\chapter{State and buffers}
\MARSChapterHeading{State and buffers}
\label{revision:state-and-buffers}

The line rasterizer made two decisions from one saved error value. A generation
of Conway's Game of Life extends that problem to an entire array: every cell's
next value depends on the same old board. If one cell reads a neighbor that has
already changed, the program computes a different process.

We will evolve a finite eight-by-eight board, render each completed generation
and retain the preceding generation in a second allocation. The complete program
is \texttt{life-buffers.asm} in \texttt{revision/labs/files}. Unlike the earlier
small examples, this chapter prints the parts that establish the new state and
ownership rules; the public file supplies the complete caller, stack-frame setup
and familiar rendering loops.

You should be able to distinguish a source snapshot from a partially built result,
reason about a padded array, preserve live values across a helper call, and identify
the boundary at which a result can become current. The independent exercise adds
a changed-cell count without changing the generation that is computed.

\section{Define a simultaneous update}

A cell is either dead, represented by byte zero, or alive, represented by byte one.
Its neighbors are the eight surrounding cells, including diagonals but excluding
itself. A dead cell becomes alive when exactly three neighbors are alive. A live
cell remains alive with two or three; all other next states are dead.\footnote{These
are Conway's Life rules as introduced in Martin Gardner's October 1970
\emph{Scientific American} column, pp.~120--123. A scan is available at
\url{https://web.stanford.edu/class/sts145/Library/life.pdf}.}

Our board is finite. Positions outside its eight-by-eight interior remain dead
and do not wrap to the opposite edge. This boundary condition is part of the
program's contract. It gives different edge behavior from an infinite board or
a board whose opposite edges are connected.

The rule has three cases. If the neighbor count is three, the
next value is one regardless of the current cell. If the count is two, retain the
current cell. Otherwise, write zero. The old cell is needed only in the second
case; the destination's previous value is never an input.

Consider three live cells at $(2,3)$, $(3,3)$ and $(4,3)$. The next generation has
live cells at $(3,2)$, $(3,3)$ and $(3,4)$. After another update, the horizontal
pattern returns. Figure~\ref{fig:revision-life-generations} shows the two distinct
states and the allocations that hold them.

\begin{figure}[htbp]
  \centering
  \begin{tikzpicture}[x=3.8mm,y=-3.8mm,alt={Three eight by eight grids. Generation zero is current in board A, with live cells two comma three, three comma three and four comma three. Generation one is current in board B, with live cells three comma two, three comma three and three comma four. Generation two is again current in board A and repeats the first horizontal pattern.}]
    \foreach \base/\title in {0/{0: board A},11/{1: board B},22/{2: board A}} {
        \node[font=\small,anchor=south] at (\base+4,-0.8) {\title};
        \foreach \n in {0,...,8} {
            \draw[black!45,thin] (\base+\n,0) -- (\base+\n,8);
            \draw[black!45,thin] (\base,\n) -- (\base+8,\n);
          }
      }
    \foreach \x/\y in {2/3,3/3,4/3,14/2,14/3,14/4,24/3,25/3,26/3} {
        \fill[black] (\x,\y) rectangle ++(1,1);
      }
    \draw[-{Stealth}] (8.4,4) -- (10.6,4);
    \draw[-{Stealth}] (19.4,4) -- (21.6,4);
  \end{tikzpicture}
  \caption{Each arrow computes a complete destination from the preceding snapshot, then exchanges the current and destination roles. The repeated final image does not mean that no update occurred.}
  \label{fig:revision-life-generations}
\end{figure}


An in-place row-major sweep changes this result. It first creates the new cell
at $(3,2)$. When it then considers $(4,2)$, the old board has two live neighbors
there, so the correct next value is zero. But the newly written $(3,2)$ would be
counted as a third neighbor, producing an incorrect birth. The rule has not changed;
the program has supplied values from two different generations to one decision.

\section{Separate storage from its current role}

Allocate two boards, called \texttt{board\_a} and \texttt{board\_b}. For one update,
treat the current board as read-only and write every next cell into the other
board. When all next cells exist, exchange the pointers that name the two roles.
No board-sized copy is needed for that exchange.

The allocation names remain fixed. After generation zero, A is current and B is
available as a destination. After one update, B is current and A retains generation
zero. After two, A is current again and B retains generation one. An allocation's
address does not say whether it is the current snapshot or the next destination;
the pointer roles and generation number establish that meaning.

Before reusing the older allocation, verify that no later decision still needs
its previous contents. In this program, the completed current board is sufficient
to calculate the next generation, so the older board can be overwritten. A different
algorithm might require more history and therefore more storage.

The update contract requires separate, nonoverlapping source and destination
boards. Giving both arguments the same address would violate the source-snapshot
requirement. It would not become safe merely because the function receives two
register arguments. A pointer is a reference to storage, not a copy of its contents.

\section{Use padding to make neighbor addresses regular}

Each board occupies 100 bytes: a ten-by-ten array containing an eight-by-eight
interior and a one-cell border. This surrounding border is called a \emph{halo}.
All 36 halo bytes start at zero and receive no program stores. An interior logical
cell $(x,y)$ is at
\[
 board+10(y+1)+(x+1),\qquad 0\leq x,y<8.
\]
The first interior byte is at offset 11 and the last at offset 88. Consecutive
interior cells differ by one byte; cells in the same column of successive rows
differ by ten. Neither stride is the framebuffer's four-byte pixel step or
32-byte row step.

For an interior cell at address $p$, the neighbor addresses have these displacements:
\[
\begin{array}{rrr}
 -11 & -10 & -9\\
 -1 & \text{center} & +1\\
 +9 & +10 & +11
\end{array}
\]
The smallest neighbor address is $board+11-11=board$. The largest is
$board+88+11=board+99$. Every neighbor load therefore lies within the 100-byte
allocation. At an edge, some loads read zero halo bytes rather than requiring a
separate bounds branch. The halo is readable state with a maintained value, not
unallocated memory that happens to follow the board.

The leaf helper below receives an interior-cell address in \texttt{\$a0} and
returns its neighbor count in \texttt{\$v0}. It may change temporary registers;
its caller must retain persistent state elsewhere. With binary cell inputs, the
result is from zero through eight. There is no load at displacement zero, because
the center is not its own neighbor.

\noindent\begin{minipage}{\linewidth}
\MARSCodeLabel{Count eight bytes from the source snapshot}
\begin{minted}{mars-idea}
count_neighbors:
    lbu $v0, -11($a0)
    lbu $t0, -10($a0)
    addu $v0, $v0, $t0
    lbu $t0, -9($a0)
    addu $v0, $v0, $t0
    lbu $t0, -1($a0)
    addu $v0, $v0, $t0
    lbu $t0, 1($a0)
    addu $v0, $v0, $t0
    lbu $t0, 9($a0)
    addu $v0, $v0, $t0
    lbu $t0, 10($a0)
    addu $v0, $v0, $t0
    lbu $t0, 11($a0)
    addu $v0, $v0, $t0
    jrc $ra
\end{minted}
\end{minipage}

\section{Build every destination cell before publishing it}

\texttt{life\_step} receives the source allocation base in \texttt{\$a0} and the
destination allocation base in \texttt{\$a1}. Both halos must be zero. The source
interior must contain only zero or one; destination interior bytes may initially
contain arbitrary values. On return, the source is unchanged, the destination
contains the complete next generation, and \texttt{\$v0} is its live-cell count.
The routine preserves saved registers and the caller's stack state.

This function calls \texttt{count\_neighbors}, so it must save its own return link.
The complete source allocates a 48-byte frame: eight-byte slots hold \texttt{\$ra}
at offset zero and \texttt{\$s0} through \texttt{\$s4} at offsets 8 through 40.
It then initializes source and destination cursors to their bases plus 11, puts
eight in \texttt{\$s2} for the remaining rows, and clears \texttt{\$s4} for the
live-cell count. \texttt{\$s3} holds the remaining columns of the current row.

Those saved registers keep the cursors, counters and partial count alive across
helper calls. In particular, the caller does not rely on \texttt{\$a0} retaining
the passed cell address. It obtains the old center from its saved source cursor
when the two-neighbor case needs it.

\noindent\begin{minipage}{\linewidth}
\MARSCodeLabel{Decide one next cell from the unchanged source}
\begin{minted}{mars-idea}
step_row:
    li $s3, 8
step_cell:
    move $a0, $s0
    balc count_neighbors
    li $t0, 3
    li $t1, 1
    beqc $v0, $t0, store_cell
    nop
    li $t0, 2
    move $t1, $zero
    bnec $v0, $t0, store_cell
    nop
    lbu $t1, 0($s0)
\end{minted}
\end{minipage}

Read the three rule paths together. A count of three reaches the store with one.
A count other than two or three reaches it with zero. A count of two loads the
old center, which is already zero or one. Every path executes the byte store,
including the path for a dead result. Skipping zero stores would allow an old
destination cell to survive incorrectly.

\noindent\begin{minipage}{\linewidth}
\MARSCodeLabel{Store the result and advance interior cursors}
\begin{minted}{mars-idea}
store_cell:
    sb $t1, 0($s1)
    addu $s4, $s4, $t1
    daddiu $s0, $s0, 1
    daddiu $s1, $s1, 1
    addiu $s3, $s3, -1
    bnezc $s3, step_cell
    nop
    daddiu $s0, $s0, 2
    daddiu $s1, $s1, 2
    addiu $s2, $s2, -1
    bnezc $s2, step_row
    nop
    move $v0, $s4
\end{minted}
\end{minipage}

After eight stores, both cursors point to the right halo of their current row.
Advancing by another two bytes skips that byte and the next row's left halo,
landing at its first interior byte. At row $r$ and column $c$, both cursors have
displacement $11+10r+c$ from their respective bases. This relation explains the
addresses and the limits of the loop. After the final row, the loop count reaches
zero before another neighbor call or store can occur.

The omitted epilogue restores the six saved slots, releases all 48 bytes and
returns through the restored link. The public source contains those instructions;
the fragment above belongs inside that complete function, after its setup.

\section{Publish a generation by exchanging roles}

The complete caller copies the immutable 64-byte seed into A's interior, skipping
the halo bytes between rows. That is generation zero, even when the requested
number of updates is zero. It keeps the current and destination bases in
\texttt{\$s0}/\texttt{\$s1}, the completed generation number in \texttt{\$s2}, and
the requested count in \texttt{\$s3}. The request is from zero through four.

The following control loop renders the current snapshot, stores its generation
number, and reaches \texttt{generation\_ready}. The renderer converts all 64
interior bytes to pixel words: one becomes \texttt{0x00ffffff}, zero becomes black.
It writes dead pixels too, so the image cannot retain a cell that died. Rendering
does not change either board.

\noindent\begin{minipage}{\linewidth}
\MARSCodeLabel{Observe a complete snapshot, then prepare its successor}
\begin{minted}{mars-idea}
render_generation:
    move $a0, $s0
    la $a1, frame
    balc render_board
    sw $s2, generation
generation_ready:
    beqc $s2, $s3, life_done
    nop
    move $a0, $s0
    move $a1, $s1
    balc life_step
generation_built:
    move $t0, $s0
    move $s0, $s1
    move $s1, $t0
    addiu $s2, $s2, 1
    bc render_generation
life_done:
    li $v0, 10
    syscall
\end{minted}
\end{minipage}

At \texttt{generation\_built}, the destination is complete but the caller has not
exchanged its pointer roles or rendered that new board. \texttt{generation\_ready}
is the later observation boundary: the current board, generation word and frame
all describe the same completed generation. This distinction lets the debugger
show a meaningful stable result while still allowing inspection inside the update.

The temporary \texttt{\$t0} retains the old current pointer during the exchange.
Two sequential moves without that temporary would lose one address. The physical
arrays remain in place; only their roles change. The old current board becomes the
next destination and retains the preceding generation until the following update
starts to overwrite it.

A fast Run may repaint only the final picture. Use the ready breakpoint to inspect
each generation; do not infer that the middle generation was skipped because the
window did not visibly display it. With the default two-step oscillator, the final
image equals the initial image, so a final-picture-only test could even accept a
program that performed no evolution at all.

\section{Test the transition, then extend its result}

Predict the default pattern through generations zero, one and two, then observe
each ready checkpoint in Debug. Record the generation word, current allocation
and other allocation's contents. After an odd number of updates B is current;
after an even number A is current. A stable two-by-two block tests a different
property: its contents stay the same although the pointer roles still exchange.

Use an edge pattern to check the boundary condition. Poison the destination
interior with nonbinary values to confirm that every result is computed and
stored. Check both halo contents and the absence of stores to them. A halo that
is temporarily overwritten and later restored can still contaminate another
neighbor count. Likewise, source preservation matters throughout an update, not
only at its final checkpoint.

The separate exercise implements \texttt{life\_changes} with the same board
contract, returning the number of interior cells whose new value differs from
their old value. That is not the live-cell count: the default oscillator has
three live cells in either orientation, but four cells change on each transition.
A stable block has four live cells and zero changes. You may count differences
while building the result or compare the complete destination with the retained
source afterward. In either design, the next board must be identical to the
ordinary Life update.

The ring-buffer companion applies related reasoning to queued records. It uses
head, tail and occupied-byte counts to decide which storage can be read or reused,
and its failed operations must preserve caller state. That lab adds wraparound
and variable-length records; complete the snapshot reasoning here before taking
on those additional interface rules.
